\chapter{Coupled observer, self-inertia and genealogical conservation of information}
\label{chap:observador-autoinercia-landauer}

\section{Joint state and observer residence}

The elementary dynamic description comprises the observed section, the
realized observer, and the record produced by both. We write
\begin{equation}
 \boxed{\Gamma=(x,\mathcal O,m)\in
 \mathcal X_{\rm TPK}\times\mathfrak O\times\mathfrak M.}
 \label{eq:rm-obs-joint-state}
\end{equation}
The coordinate \(x\) preserves cylinder, route, sheet, orientation, carry,
boundary, and possibilities of prolongation; \(\mathcal O\) specifies the
reader subsystem and its resolution; \(m\) preserves the distinguishable memory
of the coupling. The joint state constitutes the primary domain of
measurement and inertial response.

An observer at depth \(n\) is represented by a chart
\begin{equation}
 q_{\mathcal O,n}:\mathcal X_{\rm TPK}^{(n)}
 \longrightarrow Y_{\mathcal O,n}.
 \label{eq:rm-obs-observer-chart}
\end{equation}
For \(x\in\mathcal X_{\rm TPK}^{(n)}\), its observational horizon is the
fibre
\begin{equation}
 \mathfrak H_{\mathcal O,n}(x)
 =q_{\mathcal O,n}^{-1}(q_{\mathcal O,n}(x)).
 \label{eq:rm-obs-horizon-fiber}
\end{equation}
The fiber gathers the sections indistinguishable for the declared chart.
Refinement preserves naturality when there are maps
\(r_{n+1,n}\) and \(s_{n+1,n}\) such that
\begin{equation}
 q_{\mathcal O,n}\circ r_{n+1,n}
 =s_{n+1,n}\circ q_{\mathcal O,n+1}.
 \label{eq:rm-obs-observer-naturality}
\end{equation}
This identity makes finite observation compatible with deep genealogy:
the result of truncating a refined reading coincides with the
reading of the truncated state.

\begin{definition}[Realized observer]
An implemented observer is a tuple
\begin{equation}
 \mathcal O_a=(m_0,U_a,q_a,\mathcal C_a)
 \label{eq:rm-obs-realized-observer}
\end{equation}
formed by a preparation \(m_0\), a coupling \(U_a\), a pointer
reader \(q_a\), and a compatibility and prolongation condition
\(\mathcal C_a\). Its function is to produce a distinguishable record
and use it as a condition for continuation of the route.
\end{definition}

The definition situates the reading frame within the
system. The published coordinate depends on the coupling, the
resolution, and the memory; the enriched
section preserves the totality from which that coordinate
arises.

\section{Measurement context and conditioning by fibers}

Let \((\Omega_S,\Sigma_S)\) and \((\Omega_M,\Sigma_M)\) be measurable spaces
of sections of the system and the apparatus. The context \(a\) uses a
measurable coupling
\begin{equation}
 U_a:\Omega_S\times\Omega_M
 \longrightarrow\Omega_S\times\Omega_M
 \label{eq:rm-obs-measurable-coupling}
\end{equation}
and a reader \(q_a:\Omega_M\to Y_a\). The output is
\begin{equation}
 \operatorname{Out}_a(x)
 =q_a\!\left(\operatorname{pr}_M U_a(x,m_0)\right).
 \label{eq:rm-obs-output}
\end{equation}
For an outcome \(y\), the measurable fiber is defined as
\begin{equation}
 \Omega_{a,y}
 =\{x\in\Omega_a:\operatorname{Out}_a(x)=y\}.
 \label{eq:rm-obs-outcome-fiber}
\end{equation}
The Update of Genealogy Is the Restriction
\begin{equation}
 \Omega_a\longmapsto\Omega_{a,y}.
 \label{eq:rm-obs-genealogical-update}
\end{equation}
At finite depth, the operation preserves the prefixes that admit an
extension in \(\Omega_{a,y}\). The memory of the result is incorporated into the
prolongation condition; the traversed route remains recorded.

If \(\mu_a(\Omega_{a,y})>0\), the measure conditioned is
\begin{equation}
 \mu_{a,y}(B)
 =\frac{\mu_a(B\cap\Omega_{a,y})}
 {\mu_a(\Omega_{a,y})}.
 \label{eq:rm-obs-conditional-measure}
\end{equation}
This expression constitutes the probabilistic chart of the relative state.
Its spectral realization is obtained through cylindrical projections.

\section{cylindrical PVM and Lüders correspondence}

Suppose that \(\Omega_{a,y}\) is a finite union of disjoint cylinders
of a common depth \(n\). Let \(P_y\) be the sum of their
orthogonal projectors in the Hilbert realization. For a density
state \(\rho\) and a common PVM of cylindrical events, the Born probability
is
\begin{equation}
 p(y)=\Tr(\rho P_y).
 \label{eq:rm-obs-born}
\end{equation}
When \(p(y)>0\), the update by Lüders \cite{Luders1951} is
\begin{equation}
 \rho_y=\frac{P_y\rho P_y}{\Tr(\rho P_y)}.
 \label{eq:rm-obs-luders}
\end{equation}

\begin{theorem}[Equivalence between fiber restriction and Lüders]
\label{thm:rm-obs-fiber-luders}
Be \(\mathcal A_n\) the algebra of sucesos cilindricos to depth
\(n\). Supongamos that
\begin{equation}
 \mu_a(B)=\Tr(\rho P_B),
 \qquad B\in\mathcal A_n,
 \label{eq:rm-obs-measure-pvm}
\end{equation}
and that \(P_B\) and \(P_y\) belong to the same projective spectral measure.
Then
\begin{equation}
 \boxed{
 \Tr(\rho_yP_B)
 =\frac{\mu_a(B\cap\Omega_{a,y})}
 {\mu_a(\Omega_{a,y})}.}
 \label{eq:rm-obs-luders-fiber-equivalence}
\end{equation}
\end{theorem}

\begin{proof}
Membership in a common PVM gives
\(P_BP_y=P_{B\cap\Omega_{a,y}}\). Cyclicity of the trace and
\cref{eq:rm-obs-measure-pvm} yield
\[
 \Tr(\rho_yP_B)
 =\frac{\Tr(\rho P_yP_BP_y)}{\Tr(\rho P_y)}
 =\frac{\Tr(\rho P_{B\cap\Omega_{a,y}})}
 {\Tr(\rho P_y)},
\]
which coincides with \cref{eq:rm-obs-conditional-measure}.
\end{proof}

The hypothesis of a common PVM is substantive. Two mutually
unbiased contexts admit their own projective descriptions and require a noisy POVM
or a dilation for a joint reading. The correspondence of
\cref{thm:rm-obs-fiber-luders} is asserted within the declared spectral
context.

\section{Quantum instrument, dilation, and channel preservation}

Let \(\widehat U_a\) be the unitary realization of the coupling between
system and apparatus, and let \(|m_0\rangle\) be the pointer preparation. A
refined basis \(|m_{y,\alpha}\rangle\) induces the Kraus operators
\begin{equation}
 M_{y,\alpha}
 =\langle m_{y,\alpha}|\widehat U_a|m_0\rangle,
 \qquad
 \sum_{y,\alpha}M_{y,\alpha}^{\dagger}M_{y,\alpha}=I.
 \label{eq:rm-obs-kraus}
\end{equation}
The instrument and its effect are
\begin{equation}
 \mathcal I_y(\rho)
 =\sum_\alpha M_{y,\alpha}\rho M_{y,\alpha}^{\dagger},
 \qquad
 E_y=\sum_\alpha M_{y,\alpha}^{\dagger}M_{y,\alpha}.
 \label{eq:rm-obs-instrument}
\end{equation}
The probability and the relative state take the form
\begin{equation}
 p(y)=\Tr(\rho E_y),
 \qquad
 \rho_y=\frac{\mathcal I_y(\rho)}{p(y)}.
 \label{eq:rm-obs-general-conditioned-state}
\end{equation}
The non-selective channel is
\begin{equation}
 \mathcal E_a(\rho)=\sum_y\mathcal I_y(\rho).
 \label{eq:rm-obs-unconditioned-channel}
\end{equation}

\begin{proposition}[Exact Conservation of the Global Channel]
\label{prop:rm-obs-global-channel}
The isometry
\begin{equation}
 V_a:\mathcal H_S\longrightarrow\mathcal H_S\otimes\mathcal H_M,
 \qquad
 V_a|\psi\rangle=\widehat U_a(|\psi\rangle\otimes|m_0\rangle)
 \label{eq:rm-obs-stinespring-isometry}
\end{equation}
satisfies \(V_a^\dagger V_a=I\). Consequently,
\begin{equation}
 \sum_y p(y)=1,
 \qquad
 \Tr\mathcal E_a(\rho)=\Tr\rho,
 \label{eq:rm-obs-trace-preservation}
\end{equation}
and \(\mathcal E_a\) is completely positive and trace-preserving.
For an initial pure state, the joint state
\(V_a|\psi\rangle\) remains pure and normalized.
\end{proposition}

\begin{proof}
The unitarity of \(\widehat U_a\) and the normalization of
\(|m_0\rangle\) yield \(V_a^\dagger V_a=I\). The completeness identity for
\cref{eq:rm-obs-kraus} preserves the trace. The Kraus form proves
complete positivity, and an isometry maps unit vectors to
unit vectors.
\end{proof}

Preservation of the global channel means that the dilated evolution
retains all amplitudes and all records. Conditioning
selects a relative fiber for the reader that obtained \(y\). The sum of
the instruments reproduces the complete channel. This threefold distinction
separates evolution, reading, and relative state.

The Stinespring theorem \cite{Stinespring1955} guarantees a dilation that is isometric for every
completely positive channel. The identification of the auxiliary space with
the TPK memory is realized by an intertwiner that preserves route,
carry, register, and refinement. The abstract existence of the dilation
and the concrete genealogy of the register are two coordinable, typed results.

\section{State relative and formulation Everettian}

The relative-state formulation provides the conceptual antecedent
of this construction \cite{Everett1957}.

In an ideal premeasurement, suppose that the pointer states
\(\{|m_r\rangle\}_r\) are orthonormal and that \(\widehat U_a\) is an
isometry on the subspace spanned by
\(\{|r\rangle\otimes|m_0\rangle\}_r\). Then
\begin{equation}
 \widehat U_a(|r\rangle\otimes|m_0\rangle)
 =|r\rangle\otimes|m_r\rangle.
 \label{eq:rm-obs-ideal-premeasurement}
\end{equation}
For \(|\psi\rangle=\sum_r c_r|r\rangle\), linearity produces
\begin{equation}
 |\Psi'\rangle
 =\sum_r c_r|r\rangle\otimes|m_r\rangle.
 \label{eq:rm-obs-entangled-record}
\end{equation}
The record \(m_r\) defines the relative section
\begin{equation}
 \rho_r
 =\frac{Q_r|\Psi'\rangle\langle\Psi'|Q_r}
 {\Tr(|\Psi'\rangle\langle\Psi'|Q_r)},
 \qquad
 Q_r=I_S\otimes|m_r\rangle\langle m_r|.
 \label{eq:rm-obs-relative-state}
\end{equation}
In the Genealogical Space, the Same Result Corresponds to
\begin{equation}
 \Omega_{a,r}
 =\{x:\operatorname{Out}_a(x)=r\}.
 \label{eq:rm-obs-relative-fiber}
\end{equation}

\begin{theorem}[Relative state correspondence]
\label{thm:rm-obs-everett-correspondence}
Under the projective hypotheses of
\cref{thm:rm-obs-fiber-luders}, the relative state of
\cref{eq:rm-obs-relative-state} and the measure conditioned on
\(\Omega_{a,r}\) assign the same probabilities to every event of
the declared cylindrical PVM. The global isometry preserves the correlated superposition
of all records.
\end{theorem}

\begin{proof}
The first assertion is
\cref{eq:rm-obs-luders-fiber-equivalence}. The second follows from
\(V_a^\dagger V_a=I\) and the linear expression
\cref{eq:rm-obs-entangled-record}.
\end{proof}

The Everettian correspondence reaches the mathematical level of relative
states: the joint state preserves the channels, and each register determines
a conditioned fiber. HMT adds to each fiber a genealogy composed
of prefixes, carry, orientation, boundary, and compatible prolongations.
The macroscopic stability of the registers belongs to the theory of
decoherence of the apparatus; the ontology of physically autonomous branches
belongs to an additional interpretation. The preceding theorem retains its
force independently of that ontological choice.

